\documentclass[a4paper]{article}
\usepackage[polish]{babel}
\usepackage[cp1250]{inputenc}
\usepackage[T1]{fontenc}
\usepackage[dvips]{graphicx}

\pagestyle{headings}
\textwidth      15.5cm
\oddsidemargin    .1cm
\evensidemargin   .1cm

\begin{document}

\thispagestyle{empty}

\begin{minipage}{5cm}
\includegraphics[scale=0.7]{salomon.eps}
\end{minipage}
\begin{minipage}{10cm}
\begin{center}
{\Large Salomon\\
System przetwarzania wiedzy\\}
\end{center}
\end{minipage}

\vspace*{0.5cm}

\hrulefill

\vspace*{1cm}


\begin{flushright}
{\Large \emph{Iteration Description --- E1, ver 1.0} }
\end{flushright}



\begin{flushleft}
{ \Large Historia wersji }\[\]
\begin{tabular}{|p{2.5cm}||p{1.5cm}|p{6cm}|p{3.5cm}|}
  \hline
  Data & Wersja & Opis & Autor \\
  \hline \hline
  18.01.2006 & 1.0 & utworzenie & \emph{Leszek Grzanka}\\
  \hline
  \hline
\end{tabular}
\end{flushleft}

\newpage

\section{Zadania dla iteracji} 

Zaprojektowanie obs�ugi \textit{dataset�w} poprzez obecny silnik platformy. Zadaniem tej iteracji jest tak�e zapoznanie si� z przyk�adowymi algorytmami realizowanymi przez Wek�.

\section{Produkt iteracji} 

Implementacja zaprojektowanych interfejs�w.

\section{Plan test�w} 

\begin{itemize}

\item Prezentacja projektu oraz pustych interfejs�w.
\item Konsultacja z drug� grup� pracuj�c� nad Salomonem.
\item Konsultacja z klientem.

\end{itemize}

\section{Efekt wykonania iteracji}

Wprowadzono do kodu odpowiednie interfejsy z przeznaczeniem p�niejszego zaimplementowania.

Ze wzgl�du na du�o wy�szy priorytet zada� zwi�zanych z \textit{datasetami} nie wykonano cz�ci zadania zwi�zanego z Wek�.

\end{document}

